\section*{第 \ref{ch:g6:describing-motion} 章——描述运动：轨迹与速度}

\begin{solution}{exo:g6:describing-motion:1}
直线的（沿窗格向下）；圆周的；曲线的（拱形）；曲线的（
蜿蜒的 S 形线）。
\end{solution}

\begin{solution}{exo:g6:describing-motion:2}
轮毂：一条直线，与道路平行滑行。气门嘴：
一条优美的跳跃曲线——升起、拱起，
每转一圈触地一次。
\end{solution}

\begin{solution}{exo:g6:describing-motion:3}
匀速的；变速的（加速）；匀速的；变速的（减速）。
\end{solution}

\begin{solution}{exo:g6:describing-motion:4}
间隔相等：\omterm{def:g6:describing-motion:uniform}{匀速运动}。间隔逐渐增大（$1, 2, 4, 7$）：加速
运动——每个相等时段内经过的距离越来越多。
\end{solution}

\begin{solution}{exo:g6:describing-motion:5}
漂移的相机会把自身的运动也记录进去——这些点
便会把物体的路程与相机的运动混在一起。运动总是
相对于某物而言的，因此胶片必须确定它的
“参照物”：参见关于视角的注。
\end{solution}

\begin{solution}{exo:g6:describing-motion:6}
表：$30$ min $\to$ \qty{12}{km}，因此 $60$ min $\to$
\qty{24}{km}：渡轮的\omterm{def:g5:speed-distance-time:speed}{速度}为 $24$ 千米每\omterm{ex:g2:measuring-time:clock}{小时}。
\end{solution}

\begin{solution}{exo:g6:describing-motion:7}
$60$ min $\to$ \qty{14}{km}；$30$ min $\to$ \qty{7}{km}；$15$ min
$\to$ \qty{3.5}{km}。
\end{solution}

\begin{solution}{exo:g6:describing-motion:8}
散步者以 $12$ 千米每\omterm{ex:g2:measuring-time:clock}{小时}行走：不予采纳——这是散步\omterm{def:g5:speed-distance-time:speed}{速度}的
三倍。公交车：$10$ km 用时 $30$ min，\omterm{def:g5:speed-distance-time:speed}{速度}为 $20$
千米每\omterm{ex:g2:measuring-time:clock}{小时}——恰是城市公交车的\omterm{def:g5:speed-distance-time:speed}{速度}，予以采纳。火车：
$75$ km 用时 $15$ min，\omterm{def:g5:speed-distance-time:speed}{速度}为 $300$ 千米每\omterm{ex:g2:measuring-time:clock}{小时}——
已属顶尖水平，但对高速列车而言合理：予以采纳。
\end{solution}

\begin{solution}{exo:g6:describing-motion:9}
对面的乘客：完全没有运动——没有\omterm{def:g6:describing-motion:trajectory}{轨迹}，
\omterm{def:g5:speed-distance-time:speed}{速度}为零。对奶牛而言：是一条长长的\omterm{def:g6:describing-motion:trajectory}{直线轨迹}，\omterm{def:g5:speed-distance-time:speed}{速度}为火车全速。
两种说法都正确：二者都是以各自的视角描述推土机，
而运动总是相对于某物而言的。
\end{solution}

\begin{solution}{exo:g6:describing-motion:10}
乘车时间：$18 - 6 = 12$ 分钟，行驶 \qty{9}{km}。表：$12$
min $\to$ $9$ km，因此 $60$ min $\to$ $45$ km——恰好为 $45$
千米每\omterm{ex:g2:measuring-time:clock}{小时}：既不偏高，也不偏低，恰好正确。
\end{solution}

\begin{solution}{exo:g6:describing-motion:11}
用照相机拍摄狗取物的过程；每半秒暂停一次；标出
狗相对于栅栏柱的位置；\omterm{def:g6:measurement-in-science:measuring}{测量}间隔。间隔处处
相等：\omterm{def:g5:speed-distance-time:speed}{速度}稳定。间隔先扩大、后缩小（或
大小交替）：冲刺后滑行。这些点给出了
结论，无论家人怎样喊叫。
\end{solution}

\begin{solution}{exo:g6:describing-motion:12}
从房间内看：曲线——小球向前并向下拱起，从桌沿
落到地板。从以小球前进\omterm{def:g5:speed-distance-time:speed}{速度}滚动的小车上看：
小球既不领先也不落后——它
看上去是\emph{竖直向下}地落下。同一次下落，两种视角，
而较简单的视角属于运动的观察者：这是一颗
伟大思想的种子。
\end{solution}

\begin{solution}{pb:g6:describing-motion:1}
\textbf{1.} 间隔均为 \qty{50}{m}：\omterm{def:g6:describing-motion:uniform}{匀速运动}。
\textbf{2.} 总路程为 $5 \times 50 = \qty{250}{m}$。
\textbf{3.} 有轨电车从静止出发：间隔较短并逐渐增大——
前三个间隔内加速——随后转为
稳定的 \qty{50}{m} 间隔：匀速行驶。
\textbf{4.} 最后两段间隔各为 \qty{50}{m}。
\textbf{5.} $10$ s $\to$ \qty{50}{m}，因此 $60$ s $\to$
\qty{300}{m} 每分钟。
\textbf{6.} $60$ min $\to$ $60 \times 300 = \qty{18000}{m}$：
匀速行驶时的\omterm{def:g5:speed-distance-time:speed}{速度}为 $18$ 千米每\omterm{ex:g2:measuring-time:clock}{小时}。
\textbf{7.} 表：\qty{18}{km} 每 $60$ min，因此 \qty{12}{km} 用时
$40$ 分钟。
\textbf{8.} $50 - 40 = 10$ 分钟因停车和慢行路段而损失。
\textbf{9.} $50$ min $\to$ \qty{12}{km}，因此 $60$ min $\to$
$12 \times 60 \div 50 = \qty{14.4}{km}$：全程的\omterm{def:g5:speed-distance-time:speed}{速度}为
约 $14$ 千米每\omterm{ex:g2:measuring-time:clock}{小时}——海报上的 $18$ 只在
匀速路段上准确，而不是“全城畅通”。
\textbf{10.} 全程\omterm{def:g5:speed-distance-time:speed}{速度}描述的是\emph{仿佛}整个行程
是匀速的：那个稳定\omterm{def:g5:speed-distance-time:speed}{速度}，恰好能走完
同样的 \qty{12}{km}，用时同样的 $50$ 分钟。没有任何\omterm{def:g5:speed-distance-time:speed}{速度}表显示它，
但它是规划时的合理数值，因为它已经
包含了每一次停车和慢行路段。
\textbf{11.} $50 - 4 = 46$ 分钟；$12 \times 60 \div 46 \approx
\qty{15.7}{km}$ 每\omterm{ex:g2:measuring-time:clock}{小时}——可取约 $16$。
\textbf{12.} 例如：“通勤者应当按
全程\omterm{def:g5:speed-distance-time:speed}{速度}来规划出行——今天约为 $14$ 千米每\omterm{ex:g2:measuring-time:clock}{小时}——
因为它是唯一包含等待时间的数值；
匀速段的 $18$ 美化了这条线路，却漏掉了每一次停车。”
\end{solution}
